\documentclass[12pt]{article}
\usepackage{pst-all,amssymb,graphicx,marvosym,array,fancyhdr,amsmath,multirow,pdflscape,rotating}
\usepackage{color}
\usepackage[dutch]{babel}
\usepackage{cmbright}
\usepackage[T1]{fontenc}
\usepackage{wrapfig,hyperref}


\def\dfrac#1#2{\displaystyle{\frac{#1}{#2}}}
\paperheight=297mm
\voffset=-1in
\topmargin=0.5cm
\headheight=0.5cm
\headsep=0.5cm
\textheight=27cm
\footskip=0.5cm
\paperwidth=210mm
\hoffset=-1in
\oddsidemargin=2cm
\evensidemargin=2.5cm
\textwidth=17cm
\marginparsep=0cm
\marginparwidth=0cm
\setlength{\parindent}{0pt}
%\pagestyle{fancy}

\newgray{lichtgrijs}{0.90}


\def\ds#1{\displaystyle{#1}}
\def\titel#1{\vskip 1cm
            {\large {\bf{#1}}}
            \vskip 0.5cm}
\def\subtitel#1{{\bf #1\smallskip}}



\def\stelsel#1{\left\{ \begin{array}{l}
               #1 \end{array}\right.}
\def\matrix#1{\left( \begin{array}{ccc}
               #1 \end{array}\right)}
\def\matrixtwee#1{\left( \begin{array}{cc}
               #1 \end{array}\right)}
\def\matrixdrie#1{\left( \begin{array}{ccc}
               #1 \end{array}\right)}
\def\matrixvier#1{\left( \begin{array}{cccc}
               #1 \end{array}\right)}
\def\uitgebrmatrix#1{\left( \begin{array}{ccc|c}
               #1 \end{array}\right)}
\def\kolommatrix#1{\left( \begin{array}{c}
               #1 \end{array}\right)}
\newcommand{\RR}{\mbox{\msbm R}}
\def\opdr#1{\medskip

\underline{Opdracht}\begin{enumerate}
\addtocounter{enumi}#1}

\def\eopdr{\end{enumerate}}

\def\opdrn#1{\medskip

\underline{Opdrachten}\begin{enumerate}
\addtocounter{enumi}#1}

\def\binom#1#2{\left(\begin{array}{c}
#1\\#2
\end{array}\right)}


\newfont{\bbb}{msbm9}
\newcommand{\Z }{\ensuremath{\mathbb Z}}
\newcommand{\R }{\ensuremath{\mathbb R}}

\newcommand{\ekader}[1]{\psframebox{#1}}
\newcommand{\kader}[2]{\psframebox*[fillcolor=#1,border=1pt,bordercolor=black,framesep=3pt]{#2}}

\newcommand{\kaderm}[3]{\psframebox*[fillcolor=#1,border=1pt,bordercolor=black,framesep=3pt]%
                       {\parbox{#2}{#3}}}
\newcommand{\puntenklein}[1]{\hfill\kader{white}{\small \ \ \ \ \ \ \ /{#1}}}
\newcommand{\punten}[1]{\hfill\kader{white}{\Large \ \ \ \ \ \ \ /{#1}}}
\newcommand{\blijst}{\smallskip

           \begin{list} {}{\setlength{\itemsep}{0.005em}
\setlength{\leftmargin}{2em} \setlength{\topsep}{0.05em}}}
\newcommand{\elijst}{\end{list}\rm \medskip}
\newcommand{\bopm}{\begin{list}
    {\it Opmerking:}{\setlength{\topsep}{0em}
    \setlength{\labelwidth}{0em}\setlength{\itemindent}{-5.75mm}
    \setlength{\leftmargin}{7.5mm}}\small\item}
\newcommand{\eopm}{\end{list}\normalsize\par}
\newcommand{\N}{{\Bbb N}}
\begin{document}
\fancyhf{}
\pagestyle{fancy}
\fancyfoot[R]{\footnotesize \thepage}
\lhead{\footnotesize \sc Onderzoekscompetenties wiskunde 6ECWW-6GRWW-6LAWW}
\rhead{\footnotesize \sc Schooljaar 2017-2018}

%Start Hoofding%%%%%%%%%%%%%%%%%%%%%%%%%%%%%%%%%%%%%%%%%%%%%%%%%%%%%%%%%%%%%%%%%%%%%%%%%%%%

\scalebox{0.96}[1]{\begin{tabular}[t]{||c||ll||}
\hline  \hline
\raisebox{-0.5cm}{\includegraphics[scale=0.4]{logoschool.eps}} &
\begin{tabular}{ll|ll}
Ond:&... & Groepsleden:& ........ \\
& & & ........ \\
& & & .............. \\
& & & .............. \\
& & &\\
\end{tabular}
& \\
\sc Hoogstraat 30 & \raisebox{-0.3cm}{\kader{yellow}{\textbf{Verslag onderzoekscompetenties wiskunde (Onderwerp .....)}}} &  \\
\sc 9700 Oudenaarde&   & \\
\hline \hline
\end{tabular}\\[1ex]}
%Einde Hoofding%%%%%%%%%%%%%%%%%%%%%%%%%%%%%%%%%%%%%%%%%%%%%%%%%%%%%%%%%%%%%%%%%%%%%%%%%%%%
\tableofcontents


\newpage
\section{test}
Hier kan je beginnen schrijven...\cite{test}%deze laatste regel is een link naar je bronvermelding

\bigskip
\begin{wrapfigure}{r}{0.25\textwidth}
    \centering
  	\includegraphics[scale=0.3]{zonderGRM.JPG}
\end{wrapfigure}

Zo kan je dus figuren toevoegen in je bestand en je tekst eromheen ``wrappen''. Via de knop ``PROJECT'' bovenaan zie je welke files er in dit document staan.

\subsection{test}

En zo typ je wiskunde:
\begin{itemize}
\item Stelsels:
$$\left\{\begin{array}{rcl}
f(x)&=& 6\sin\left[\dfrac{\pi}{300}(x-323)\right]+36\\
g(x)&=& \dfrac{36}{323}x
\end{array}\right.$$

\item Matrices:
$\bordermatrix{\text{ }&K_1&K_2&K_3&K_4\cr
                L_a& 1 &  0  & 0 & 1\cr
                L_b& 0  &  1 & 1 & 0\cr
                L_c& 1 & 0 & 1 & 0\cr
                L_d& 0 & 1 & 0 & 1}$ of
$\begin{pmatrix}
1 & 0 & 0 & | & 1 & 0\\
0 & 1 & 0 & | & 1 & 0\\
1 & 0 & 1 & | & 0 & 0\\
- & - & - & | & - & -\\
1 & 0 & 0 & | & 0 & 1\\
0 & 1 & 0 & | & 0 & 1
\end{pmatrix}$

\item Besluiten in fancy kadertjes:

\kaderm{yellow}{17cm}{De sinusfunctie $g(x)$ met bovenstaande waarden heeft als voorschrift:\begin{center}$g(x)= 7\sin\left[\dfrac{\pi}{19}(x-9)\right]+7$\end{center}}
\end{itemize}

\begin{thebibliography}{1}
% hier komen je bronnen. 
\bibitem{test}
\url{http://www.google.be}



\end{thebibliography}


\end{document}